
\documentclass{article}
\usepackage{colortbl}
\usepackage{makecell}
\usepackage{multirow}
\usepackage{supertabular}

\begin{document}

\newcounter{utterance}

\centering \large Interaction Transcript for game `codenames', experiment `opponent\_easy', episode 4 with reflexion{-}Qwen3.5{-}27B.
\vspace{24pt}

{ \footnotesize  \setcounter{utterance}{1}
\setlength{\tabcolsep}{0pt}
\begin{supertabular}{c@{$\;$}|p{.15\linewidth}@{}p{.15\linewidth}p{.15\linewidth}p{.15\linewidth}p{.15\linewidth}p{.15\linewidth}}
    \# & $\;$A & \multicolumn{4}{c}{Game Master} & $\;\:$B\\
    \hline

    \theutterance \stepcounter{utterance}  
    & & & \multicolumn{2}{p{0.3\linewidth}}{
        \cellcolor[rgb]{0.9,0.9,0.9}{
            \makecell[{{p{\linewidth}}}]{
                \texttt{\tiny{[GM$|$GM]}}
                {'hidden': {'team': ['pumpkin', 'pound', 'tooth', 'fire', 'plot', 'kangaroo', 'bow', 'ketchup'], 'innocent': ['fan', 'capital', 'ivory', 'orange', 'grass', 'slip', 'africa'], 'opponent': [], 'assassin': ['antarctica']}, 'revealed': {'team': {'team': [], 'innocent': [], 'opponent': [], 'assassin': []}, 'opponent': {'team': [], 'innocent': [], 'opponent': [], 'assassin': []}}}
            }
        }
    }
    & & \\ \\

    \theutterance \stepcounter{utterance}  
    & & \multicolumn{4}{p{0.6\linewidth}}{
        \cellcolor[rgb]{0.9,0.9,0.9}{
            \makecell[{{p{\linewidth}}}]{
                \texttt{\tiny{[P1$\langle$GM]}}
                \texttt{You have attempted this same task before and did not succeed. These are the plans you wrote after those attempts.} \\
\texttt{Trial \#1: New plan: My previous strategy failed because I provided clues that were too broad or ambiguous, leading my partner to guess opponent words (e.g., "red" triggered "heart", "movie" triggered "film", "game" triggered "space"). I also listed opponent words as targets in my third turn ("film"), which was a critical error. Next time, I will strictly verify that my clue is semantically unique to my remaining team words and has no strong association with any remaining opponent, distractor, or kill words. I will also ensure my TARGETS list only contains words from my own team. For the remaining words (pumpkin, pound, tooth, fire, plot, kangaroo, bow, ketchup), I will use more specific, less common associations (e.g., "Halloween" for pumpkin, "currency" for pound, "enamel" for tooth) to minimize the risk of triggering opponent words like "pistol" or "soul". I will prioritize clues that link multiple safe words only if the association is unambiguous.} \\
\\ 
\texttt{Trial \#2: New plan: My previous attempt failed because I selected a clue ("HALLOWEEN") that was too narrow, covering only two words while leaving six unaddressed, and I failed to identify a unifying theme for the entire set. I also did not account for the risk of the clue triggering distractors (e.g., "orange" for pumpkin). Next time, I will first scan my team's words for a strong, common semantic thread that covers the maximum number of items (e.g., "pound" as a unit of weight, "pound" as a currency, "pound" as a verb, or perhaps a theme like "things that can be struck" or "parts of a house" if applicable, though here "pound" is the outlier). If no single word covers all, I will prioritize a clue that covers the largest safe cluster (e.g., "pumpkin, ketchup, fire" might relate to "orange" or "red" but those are distractors, so I must avoid them). I will specifically look for a clue like "STRIKE" (for pound, bow, fire, plot?) or "WEIGHT" (for pound, kangaroo?) and rigorously check that the clue does not strongly associate with any distractor (fan, capital, ivory, orange, grass, slip, africa) or the kill word (antarctica). I will ensure the TARGETS list includes every word the clue reasonably applies to, maximizing the count while maintaining safety.} \\
\\ 
\texttt{Trial \#3: New plan: My previous attempt was aborted because I violated the explicit rule stating the clue "cannot be one of the words in the lists." I used "ORANGE" as a clue, but "orange" is explicitly listed as a distractor word. To fix this, I must strictly cross{-}reference my chosen clue against the full list of team words, opponent words, distractors, and kill words before submitting. For the current set (pumpkin, pound, tooth, fire, plot, kangaroo, bow, ketchup), I will select a clue that is semantically related to a subset of these words but is not present in any forbidden list. For example, "SQUASH" relates to pumpkin and ketchup (as a condiment or vegetable) and is not in the forbidden lists, or "RED" (if not a distractor, though "orange" is, so I must be careful with colors) or "HARVEST" for pumpkin. I will ensure the TARGETS list only contains words from my team that the clue actually describes, and I will double{-}check that the clue word itself is not "orange", "fan", "capital", "ivory", "grass", "slip", "africa", or "antarctica".} \\
\\ 
\\ 
\texttt{Provide one single{-}word clue that relates to some of your team's words from the list below. You can choose to how many and to which words your clue relates to.} \\
\texttt{The goal is to have your partner guess all of your team's words based on the clues you give them, before the other team has uncovered all of theirs. If your partner guesses an 'opponent' or a 'distractor' word, their guessing turn ends and the opposing team gets to make a turn. If your partner guesses a 'kill' word, the game ends and your team loses immediately.} \\
\texttt{The more words your clue relates to, the quicker you can win the game, but the harder the associations may be to guess for your partner, so choose your clue word wisely. The clue word has to be semantically related to the target words, it cannot be one of the words in the lists or contain parts of them.} \\
\\ 
\texttt{Always give your single{-}word clue and your comma{-}separated list of related target words in the following format and make your answers as short as possible, never include any other text than is required in this form:} \\
\\ 
\texttt{CLUE: <WORD>} \\
\texttt{TARGETS: <TARGETS>} \\
\\ 
\texttt{Your team words are: pumpkin, pound, tooth, fire, plot, kangaroo, bow, ketchup.} \\
\texttt{Your opponent's team words are: .} \\
\texttt{Distractor words are: fan, capital, ivory, orange, grass, slip, africa.} \\
\texttt{The kill words are: antarctica.} \\
            }
        }
    }
    & \\ \\

    \theutterance \stepcounter{utterance}  
    & \multicolumn{4}{p{0.6\linewidth}}{
        \cellcolor[rgb]{0.9,0.9,0.9}{
            \makecell[{{p{\linewidth}}}]{
                \texttt{\tiny{[P1$\rangle$GM]}}
                \texttt{CLUE: SQUASH} \\
\texttt{TARGETS: pumpkin, ketchup} \\
            }
        }
    }
    & & \\ \\

    \theutterance \stepcounter{utterance}  
    & & & \multicolumn{2}{p{0.3\linewidth}}{
        \cellcolor[rgb]{0.9,0.9,0.9}{
            \makecell[{{p{\linewidth}}}]{
                \texttt{\tiny{[GM$|$GM]}}
                \texttt{squash} \\
            }
        }
    }
    & & \\ \\

    \theutterance \stepcounter{utterance}  
    & & & \multicolumn{2}{p{0.3\linewidth}}{
        \cellcolor[rgb]{0.9,0.9,0.9}{
            \makecell[{{p{\linewidth}}}]{
                \texttt{\tiny{[GM$|$GM]}}
                ['pumpkin', 'ketchup']
            }
        }
    }
    & & \\ \\

    \theutterance \stepcounter{utterance}  
    & & & \multicolumn{2}{p{0.3\linewidth}}{
        \cellcolor[rgb]{0.9,0.9,0.9}{
            \makecell[{{p{\linewidth}}}]{
                \texttt{\tiny{[GM$|$GM]}}
                {'word': 'pumpkin', 'assignment': 'team'}
            }
        }
    }
    & & \\ \\

    \theutterance \stepcounter{utterance}  
    & & & \multicolumn{2}{p{0.3\linewidth}}{
        \cellcolor[rgb]{0.9,0.9,0.9}{
            \makecell[{{p{\linewidth}}}]{
                \texttt{\tiny{[GM$|$GM]}}
                {'word': 'ketchup', 'assignment': 'team'}
            }
        }
    }
    & & \\ \\

    \theutterance \stepcounter{utterance}  
    & & & \multicolumn{2}{p{0.3\linewidth}}{
        \cellcolor[rgb]{0.9,0.9,0.9}{
            \makecell[{{p{\linewidth}}}]{
                \texttt{\tiny{[GM$|$GM]}}
                \texttt{opponent has won} \\
            }
        }
    }
    & & \\ \\

    \theutterance \stepcounter{utterance}  
    & \multicolumn{4}{p{0.6\linewidth}}{
        \cellcolor[rgb]{0.95,0.95,0.95}{
            \makecell[{{p{\linewidth}}}]{
                \texttt{\tiny{[P2$\rightarrow$P1]}}
                \texttt{Post{-}episode Reflexion (not sent to GM):} \\
\texttt{New plan: My previous attempt failed because I prioritized safety over progress, selecting a clue ("SQUASH") that only covered two words (pumpkin, ketchup) while leaving six words unaddressed. In a game where the goal is to guess *all* words, a low{-}yield clue is a strategic failure even if it is safe. Next time, I will first scan my team's words for a unifying theme that covers the maximum number of items (e.g., "Halloween" for pumpkin, "currency" for pound, "weapon" for bow/fire/plot, or "condiment" for ketchup). If a single clue cannot cover all words, I will prioritize the largest possible cluster of words that can be safely linked, rather than settling for a tiny, safe subset. I will rigorously verify that the chosen clue does not trigger any distractors (fan, capital, ivory, orange, grass, slip, africa) or the kill word (antarctica), but I will not sacrifice the number of targets covered unless the risk of a kill/distractor is high. For this specific set, I will look for a broader connection like "HARVEST" (pumpkin, ketchup, maybe kangaroo if thinking of food?) or "STRIKE" (pound, bow, fire, plot) to maximize the TARGETS list length.} \\
            }
        }
    }
    & & \\ \\

\end{supertabular}
}

\end{document}
